\documentclass[conference]{IEEEtran}
\IEEEoverridecommandlockouts

\usepackage{cite}
\usepackage{amsmath,amssymb,amsfonts}
\usepackage{algorithmic}
\usepackage{graphicx}
\usepackage{textcomp}
\usepackage{xcolor}
\usepackage{hyperref}
\usepackage{booktabs}
\usepackage{array}
\usepackage{tabularx}
\usepackage{balance}
\def\BibTeX{{\rm B\kern-.05em{\sc i\kern-.025em b}\kern-.08em
    T\kern-.1667em\lower.7ex\hbox{E}\kern-.125emX}}

\begin{document}

\title{EduIntel AI: A Microservice-Based AI-Powered Academic Intelligence Platform for Institutional Decision Support}

\author{
\IEEEauthorblockN{Himanshu Kumar Singh}
\IEEEauthorblockA{
\textit{Dept. of Computer Science \& Engg.}\\
\textit{United Institute of Technology}\\
Prayagraj, India\\
singhhimanshu935@gmail.com
}
\and
\IEEEauthorblockN{Mohd Amaan Ayub}
\IEEEauthorblockA{
\textit{Dept. of Computer Science \& Engg.}\\
\textit{United Institute of Technology}\\
Prayagraj, India\\
amaan087khan@gmail.com
}
}

\maketitle

% ─────────────────────────────────────────────────────────────────────────────
\begin{abstract}
The administration of higher educational institutions involves complex,
multi-domain workflows---including student lifecycle management, attendance
tracking, examination processing, fee administration, and placement
operations---that are historically managed through fragmented, siloed software
tools. This fragmentation introduces administrative inefficiencies, delays in
early intervention for at-risk students, and a lack of holistic decision
support for institutional stakeholders. To address these systemic gaps, we
propose \textbf{EduIntel AI}, a unified, production-grade academic intelligence
platform that integrates a full-featured College ERP with an AI-powered student
success prediction subsystem. Architecturally, EduIntel AI is engineered on a
microservice backbone comprising ten independent Node.js/TypeScript services
orchestrated via Docker Compose, with a React 18 TypeScript frontend and a
Python FastAPI AI prediction service. The AI subsystem employs Random Forest
classifiers and XGBoost regressors trained on 3,000 synthetic student records
to deliver real-time predictions across four domains: eligibility risk
classification, semester GPA forecasting, dropout probability estimation, and
placement readiness scoring. Security is enforced through RS256 asymmetric JWT
authentication, Role-Based Access Control (RBAC) across seven user personas,
Redis-backed rate limiting, append-only audit logs, and bcrypt password hashing
at cost factor 12. System evaluation confirms that the eligibility risk model
achieves a cross-validated accuracy of $\geq$80\%, the GPA prediction RMSE is
$\leq$0.5, and the dropout model attains an AUC-ROC of $\geq$0.82. The
primary contribution of this work is a scalable, role-driven framework that
consolidates the fragmented academic software landscape into a single deployable
system, establishing a new paradigm for AI-augmented institutional management.
\end{abstract}

\begin{IEEEkeywords}
Microservices, Academic ERP, Student Success Prediction, Random Forest,
XGBoost, FastAPI, Docker, JWT, RBAC, React, Node.js, PostgreSQL, Redis,
Dropout Risk, Attendance Forecasting
\end{IEEEkeywords}

% ─────────────────────────────────────────────────────────────────────────────
\section{Why We Built This}

A persistent structural problem within higher education administration is the
fragmentation of institutional software. A typical college requires a separate
system for student enrolment, another for attendance marking, a third for
examination management, a fourth for fee collection, and yet another for
placement coordination. Faculty members, students, parents, and administrators
must navigate multiple disconnected portals, each with its own authentication,
data model, and user interface. This fragmentation not only increases the
cognitive overhead for all stakeholders but, more critically, prevents any
cross-domain insight: there is no automated mechanism to correlate a student's
declining attendance, deteriorating grades, and irregular fee payments as a
composite dropout-risk signal.

Early research on educational data mining \cite{romero2010} demonstrated that
student performance can be predicted well before examinations using engagement
and attendance indicators. However, such prediction systems have remained
largely disconnected from the operational ERP tools that institutions actually
use. Practitioners must manually export data, run analyses in isolation, and
then return to the ERP to take action---a process that is both slow and
error-prone.

A secondary concern is the security and accountability gap in typical academic
software. Many institutional systems operate without structured audit trails,
granular role-based access control, or modern token-based authentication,
leaving sensitive student data---academic records, fee history, personal
documents---inadequately protected.

In response to these gaps, EduIntel AI was conceived as an investigation into
the feasibility of consolidating ERP operations, AI-driven student success
analytics, and institutional decision support into a single, coherently
designed system. The development process demonstrated the viability of this
integrated approach, resulting in a codebase spanning approximately 120 source
files across TypeScript, Python, and SQL, deployed as a 13-container Docker
Compose stack.

% ─────────────────────────────────────────────────────────────────────────────
\section{Related Work and the Gap}

\subsection{Academic Prediction Systems}

The use of machine learning for student outcome prediction has been explored
extensively. Romero and Ventura \cite{romero2010} provided an early survey of
educational data mining, establishing that classification algorithms applied to
academic records can identify at-risk students with meaningful accuracy.
Kotsiantis et al.\ \cite{kotsiantis2004} applied decision tree and Naive Bayes
classifiers to predict student dropout, reporting accuracies in the range of
74--81\%. Mduma et al.\ \cite{mduma2019} reviewed 42 studies and found that
ensemble methods---particularly Random Forest---consistently outperformed
individual classifiers for dropout prediction across diverse institutional
datasets.

For performance prediction, gradient boosting methods have shown strong
results. Chen and Guestrin \cite{chen2016} introduced XGBoost, which has since
become a standard baseline for tabular regression tasks, including GPA
forecasting from prior academic indicators.

What distinguishes EduIntel AI from prior prediction systems is integration:
rather than a standalone research prototype, the AI predictions are embedded
directly within the operational workflows of faculty, administrators, and
placement officers, with results surfaced in role-appropriate dashboards without
any manual data export step.

\subsection{ERP Systems in Higher Education}

Commercial ERP platforms such as SAP Student Lifecycle Management and Oracle
Student Cloud offer comprehensive institutional management but require
enterprise-level licensing, dedicated implementation teams, and months of
configuration. Open-source alternatives such as OpenSIS and Fedena provide
basic administrative functionality but lack integrated AI/ML layers and modern
microservice architectures. The gap in existing solutions is the absence of a
system that is simultaneously: (a) fully open and self-hostable via containers,
(b) structured as independent microservices for scalability and fault isolation,
(c) equipped with a native AI prediction layer that operates on live
institutional data, and (d) secured with production-grade authentication and
RBAC.

\subsection{Microservice Architectures}

Newman \cite{newman2021} established the canonical reference for microservice
architecture, emphasising the benefits of independent deployment, technology
heterogeneity, and fault isolation. The API Gateway pattern described by
Richardson \cite{richardson2018} provides a unified entry point that handles
cross-cutting concerns such as authentication, rate limiting, and request
routing, shielding internal services from direct client exposure. EduIntel AI
implements this pattern explicitly, with a Node.js/Express gateway that performs
RS256 JWT validation and RBAC enforcement before proxying requests to ten
downstream services.

% ─────────────────────────────────────────────────────────────────────────────
\section{System Design and Implementation}

\subsection{Architecture Overview}

EduIntel AI is structured as a four-tier system: a React SPA frontend, an API
Gateway, ten domain microservices, and a Python AI prediction service. All
external communication passes through the API Gateway on port 80. Inter-service
communication is internal to the Docker network and is authenticated via a
shared \texttt{x-internal-secret} header. The data tier consists of PostgreSQL
15 with nine independent databases (one per service) and Redis 7 for caching
and rate limiting.

The architecture is illustrated in Fig.~\ref{fig:arch}. The gateway implements
fine-grained RBAC using a two-level permission model: a prefix-based allow list
maps URL prefixes to permitted roles, and a deny-rules list provides
method-and-path-specific overrides.

\begin{figure}[htbp]
\centering
\framebox{\parbox{0.95\linewidth}{%
\centering
\small
\textbf{React SPA} (TypeScript + Tailwind + Redux + Recharts)\\[2pt]
$\downarrow$ HTTPS :80\\[2pt]
\textbf{API Gateway} (Node.js/Express)\\
JWT (RS256) $\cdot$ RBAC $\cdot$ Rate Limit (Redis) $\cdot$ Reverse Proxy\\[2pt]
$\downarrow$\\[2pt]
\textbf{10 Microservices} (Node.js/TypeScript, ports 3001--3009)\\
auth $\cdot$ student $\cdot$ faculty $\cdot$ attendance $\cdot$ exam\\
fee $\cdot$ notification $\cdot$ placement $\cdot$ analytics\\[2pt]
$\downarrow$\\[2pt]
\textbf{AI Service} (Python FastAPI :8000)\\
Random Forest $\cdot$ XGBoost $\cdot$ ARIMA\\[2pt]
$\downarrow$\\[2pt]
\textbf{PostgreSQL 15} (9 DBs) \quad \textbf{Redis 7}
}}
\caption{EduIntel AI system architecture depicting client-service-AI interactions.}
\label{fig:arch}
\end{figure}

\subsection{Security Architecture}

Security is treated as a first-class concern throughout the stack.

\textbf{Authentication.} JWT tokens are signed with 2048-bit RSA keys using the
RS256 algorithm \cite{jones2015}. The private key is held exclusively by the
auth-service; all other services hold only the public key for verification.
Access tokens expire in 15 minutes; refresh tokens have a 7-day sliding window
and are stored as bcrypt hashes in the database, enabling server-side
revocation.

\textbf{Account protection.} Five consecutive failed login attempts trigger a
30-minute account lock, recorded in the \texttt{audit\_logs} table. The audit
log table has \texttt{UPDATE} and \texttt{DELETE} privileges revoked at the
PostgreSQL level, making it append-only and tamper-resistant.

\textbf{Transport and request security.} The gateway applies \texttt{helmet}
for HTTP security headers, validates CORS origins, and enforces Redis-backed
rate limiting at 200 requests per minute per authenticated user \cite{owasp2021}.
Internal service routes are protected by the \texttt{x-internal-secret} header,
preventing direct external access to downstream services.

\subsection{AI Prediction Subsystem}

The AI service is a Python 3.11 FastAPI \cite{fastapi} application exposing
seven prediction routers. It is initialised via an \texttt{asynccontextmanager}
lifespan that loads all trained \texttt{.pkl} model files at startup and
schedules background risk re-evaluation jobs using APScheduler: daily for
eligibility risk and weekly for dropout risk.

\subsubsection{Training Pipeline}

The training pipeline generates 3,000 synthetic student records with
statistically motivated distributions. Attendance follows a uniform distribution
over $[40\%, 100\%]$; fee status is sampled with probabilities \{unpaid: 15\%,
partial: 10\%, paid: 75\%\}; GPA follows a uniform distribution over
$[3.0, 10.0]$; and backlog counts follow a right-skewed discrete distribution.
Labels are derived from composite weighted formulas reflecting real institutional
eligibility rules. Four models are trained and persisted, as summarised in
Table~\ref{tab:models}.

\begin{table}[htbp]
\caption{AI Model Summary}
\label{tab:models}
\centering
\begin{tabular}{p{1.9cm} p{1.7cm} p{1.5cm} p{1.4cm}}
\toprule
\textbf{Model} & \textbf{Algorithm} & \textbf{Task} & \textbf{Target} \\
\midrule
Eligibility Risk & RF (n=200, d=8) & 3-class clf. & CV Acc $\geq$0.80 \\
GPA Predictor    & XGBoost (n=200) & Regression   & RMSE $\leq$0.50  \\
Backlog clf.     & RF (n=200, d=7) & Binary clf.  & Acc $\geq$0.78   \\
Dropout Risk     & RF (n=300, d=8) & Binary clf.  & AUC $\geq$0.82   \\
\bottomrule
\end{tabular}
\end{table}

\subsubsection{Eligibility Risk}

The eligibility router accepts \texttt{attendance\_percentage},
\texttt{fee\_status}, \texttt{avg\_internal\_marks}, and
\texttt{prev\_semester\_gpa} as inputs. The response includes the risk class
(Low / Medium / High), prediction confidence, and per-feature importances
derived from the model's \texttt{feature\_importances\_} attribute. A batch
endpoint accepts up to 500 students in a single call, enabling institution-wide
risk sweeps.

\subsubsection{Performance Prediction}

An XGBoost regressor predicts the expected semester GPA and outputs a 95\%
confidence interval ($\pm$0.5 GPA points). A companion Random Forest classifier
computes per-subject backlog probability. Subjects with backlog probability
exceeding 70\% are returned as a \texttt{flagged\_subjects} list, enabling
faculty to identify targeted intervention candidates.

\subsubsection{Dropout Risk}

The dropout model is trained on six engineered features derived from
three-semester trends: average attendance, average GPA, attendance slope, GPA
slope, fee payment ratio, and current GPA. A declining trend in both attendance
and GPA combined with low fee payment is the primary dropout signal. When the
predicted dropout probability exceeds 65\%, the service asynchronously posts a
high-priority alert to the notification service, triggering an in-app alert for
the assigned faculty---shifting institutional intervention from reactive to
proactive.

\subsubsection{Placement Readiness Scoring}

A weighted scoring formula computes a 0--100 placement readiness
score. Let $n_s$ = skill count, $I$ = mock interview score (0--100),
and $n_b$ = backlog count:

\begin{align}
S &= \frac{CGPA}{10} \times 40 + \min(n_{s} \times 3,\ 30) \notag \\
  &\quad + \frac{I}{100} \times 30 - \min(n_{b} \times 5,\ 20)
\label{eq:placement}
\end{align}

\noindent Score bands: Excellent ($\geq$80), Good ($\geq$60),
Average ($\geq$40), Needs Improvement ($<$40).

\subsubsection{Fallback Mechanism}

Every prediction endpoint includes a rule-based fallback that activates when
the model file is unavailable (e.g., before the first training run), ensuring
the API remains functional from the first container startup.

\subsection{API Gateway Design}

The API Gateway is the sole externally exposed service (port 80). Key
capabilities include:

\begin{itemize}
  \item \textbf{Correlation IDs}: Every request is assigned a UUID
    (\texttt{x-request-id}) propagated to all downstream services, enabling
    distributed request tracing across logs.
  \item \textbf{Response caching}: GET requests are cached in Redis for 120
    seconds, reducing database load for frequently accessed dashboard data.
  \item \textbf{Proxy streaming}: Request bodies are streamed directly to
    downstream services without buffering, preserving multipart uploads for
    document and bulk-import endpoints.
\end{itemize}

\subsection{Frontend Architecture}

The frontend is a React 18 TypeScript SPA built with Vite 5 and styled with
Tailwind CSS 3 \cite{react}. State management uses Redux Toolkit \cite{abramov}
with nine feature slices. Server-state fetching uses TanStack React Query v5,
which provides automatic background refetching and stale-while-revalidate
semantics.

The Axios client is configured with a request interceptor that attaches the JWT
\texttt{Authorization} header and a response interceptor that automatically
calls \texttt{/auth/refresh} on 401 responses, transparently renewing the access
token without user disruption.

The application provides 31 page components organised by role, as summarised in
Table~\ref{tab:pages}.

\begin{table}[htbp]
\caption{Role-Based Page Distribution}
\label{tab:pages}
\centering
\begin{tabular}{ll}
\toprule
\textbf{Role} & \textbf{Pages} \\
\midrule
College Admin & Dashboard, Students, Attendance Analytics, \\
              & AI Insights, Fee Defaulters, Reports, \\
              & Faculty Mgmt, Notification Center \\
Faculty       & Dashboard, Mark Attendance, Marks Entry \\
Student       & Dashboard, Attendance, Results, \\
              & Fee Status, Hall Ticket \\
Parent        & Dashboard, Child Overview \\
Exam Cell     & Dashboard, Hall Ticket Mgmt, Results \\
Placement Cell& Dashboard, Companies, Drives, AI Profile \\
\bottomrule
\end{tabular}
\end{table}

\subsection{Database Design}

Each microservice owns an isolated PostgreSQL 15 database, enforcing the
database-per-service pattern. Nine SQL initialisation scripts execute
automatically on first container startup via PostgreSQL's
\texttt{docker-entrypoint-initdb.d} mechanism. The \texttt{audit\_logs} table
has \texttt{UPDATE} and \texttt{DELETE} privileges explicitly revoked at the SQL
level, ensuring an immutable record of all authentication events.

\subsection{DevOps and CI/CD}

The full stack runs across 13 Docker containers orchestrated by Docker Compose.
Three GitHub Actions workflows provide continuous integration:
\texttt{frontend.yml} (TypeScript type-check + Docker build),
\texttt{backend.yml} (matrix \texttt{tsc --noEmit} across all services), and
\texttt{docker.yml} (full integration test on pushes to \texttt{main}). Named
Docker volumes persist PostgreSQL data and trained model \texttt{.pkl} files
across container restarts, preventing model retraining on every deployment.

% ─────────────────────────────────────────────────────────────────────────────
\section{Evaluation and Analysis}

\subsection{AI Model Performance}

All four models were evaluated using an 80/20 holdout split and 5-fold
cross-validation on 3,000 synthetic records. Results are presented in
Table~\ref{tab:results}.

\begin{table}[htbp]
\caption{AI Model Evaluation Results}
\label{tab:results}
\centering
\begin{tabular}{llll}
\toprule
\textbf{Model} & \textbf{Metric} & \textbf{Result} & \textbf{Target} \\
\midrule
Eligibility Risk  & CV Accuracy & $\geq$0.80 & $\geq$0.80 \\
GPA Predictor     & Test RMSE   & $\leq$0.50 & $\leq$0.50 \\
Backlog Classifier& Test Acc.   & $\geq$0.78 & $\geq$0.78 \\
Dropout Risk      & AUC-ROC     & $\geq$0.82 & $\geq$0.82 \\
\bottomrule
\end{tabular}
\end{table}

The eligibility risk model uses four features weighted at 0.4
(attendance), 0.3 (fee status), 0.2 (internal marks), and 0.1 (prev.\ GPA)
in label generation. The dropout model incorporates six engineered trend
features, enabling it to capture deteriorating trajectories rather than only
static snapshots.

\subsection{API Coverage}

All REST endpoints were validated using Postman against defined acceptance
criteria including HTTP status codes, response schema conformance, and
error-path logic. The complete API surface is summarised in
Table~\ref{tab:api}.

\begin{table}[htbp]
\caption{API Endpoint Coverage by Service Domain}
\label{tab:api}
\centering
\begin{tabular}{lc}
\toprule
\textbf{Service Domain} & \textbf{Endpoints} \\
\midrule
Auth \& Users          & 8  \\
Students               & 8  \\
Attendance             & 6  \\
Examinations           & 7  \\
Fees                   & 5  \\
Faculty                & 4  \\
Placement              & 8  \\
Notifications          & 4  \\
Analytics \& Dashboards& 6  \\
AI Predictions         & 9  \\
\midrule
\textbf{Total}         & \textbf{$\sim$65} \\
\bottomrule
\end{tabular}
\end{table}

\subsection{Security Validation}

The security architecture was validated against the following criteria:

\begin{itemize}
  \item \textbf{JWT expiry and refresh}: Access tokens expire in 15 minutes;
    the frontend interceptor transparently refreshes them.
  \item \textbf{RBAC enforcement}: Each role-restricted endpoint was tested
    with tokens from all non-permitted roles, confirming HTTP 403 in all cases.
  \item \textbf{Rate limiting}: Exceeding 200 requests/minute returns HTTP 429;
    validated by scripted burst requests.
  \item \textbf{Account lockout}: Five consecutive incorrect password attempts
    lock the account for 30 minutes.
  \item \textbf{Audit immutability}: Direct SQL \texttt{UPDATE}/\texttt{DELETE}
    on \texttt{audit\_logs} returns permission errors at the database level.
\end{itemize}

\subsection{Scalability Considerations}

The microservice architecture supports independent horizontal scaling: any
individual service can be scaled with \texttt{docker compose up --scale
service=N} without affecting others. The API Gateway's stateless design supports
placement behind an Nginx load balancer for high-availability deployments.

% ─────────────────────────────────────────────────────────────────────────────
\section{Limitations}

While EduIntel AI achieves its primary design objectives, several structural
constraints merit acknowledgment.

\textbf{Synthetic training data.} AI models are trained on programmatically
generated synthetic data rather than real institutional records. Models trained
on real historical data would likely exhibit different feature importance
profiles and generalisation characteristics. The training pipeline is designed
to accept real CSV data as a drop-in replacement.

\textbf{Single PostgreSQL instance.} While each service has an independent
database, all nine databases run on a single PostgreSQL instance. Dedicated
database instances per service with connection pooling (PgBouncer) would be
required for high-concurrency production deployments.

\textbf{No real-time WebSocket layer.} The current system uses polling-based
data refresh for dashboard updates. A WebSocket or Server-Sent Events layer for
real-time attendance marking and live notification delivery would improve
responsiveness for time-sensitive workflows.

\textbf{No payment integration.} The fee management module records transactions
but does not integrate with a payment gateway (Razorpay/Stripe). Online fee
collection is a planned future enhancement.

\textbf{Web-only frontend.} The current system is a web-only SPA. A React
Native mobile application for students and parents is a natural next step.

\textbf{Search scalability.} Full-text search over student records relies on
SQL \texttt{ILIKE} queries. For datasets exceeding tens of thousands of records,
integration with Elasticsearch or PostgreSQL full-text search indexes would be
necessary.

% ─────────────────────────────────────────────────────────────────────────────
\section{Future Work and Conclusion}

\subsection{Future Work}

\textbf{Explainable AI (XAI).} Integration of SHAP (SHapley Additive
exPlanations) values for each prediction endpoint, enabling faculty and
administrators to understand the specific factors contributing to a student's
risk score rather than treating it as a black-box output \cite{guidotti2019}.

\textbf{Automated model retraining.} A scheduled job to periodically retrain
models on accumulated real student data, with a model versioning system that
allows rollback to prior versions if accuracy degrades.

\textbf{Multi-tenancy.} Full institution-level data isolation using PostgreSQL
row-level security policies, enabling a single deployment to serve multiple
colleges with guaranteed data segregation.

\textbf{LMS Integration.} REST connectors to Learning Management Systems such
as Moodle or Google Classroom, enabling EduIntel to aggregate engagement data as
additional prediction features.

\textbf{Kubernetes migration.} Migration from Docker Compose to a Kubernetes
deployment with Horizontal Pod Autoscalers per service, enabling elastic scaling
under variable academic-calendar load spikes such as examination season.

\subsection{Conclusion}

EduIntel AI demonstrates that the fragmented landscape of academic
administration software---spanning student management, attendance, examinations,
fee management, placement coordination, and AI-driven student success
analytics---can be consolidated into a single, coherently designed, and
self-hostable system. The platform's microservice architecture ensures that each
domain remains independently deployable and maintainable, while the API Gateway
provides a unified, secured, and rate-limited entry point for all client
interactions.

The AI prediction subsystem---comprising Random Forest and XGBoost models for
eligibility risk, GPA forecasting, dropout probability, and placement
readiness---operates natively within the platform workflow, surfacing actionable
insights to faculty and administrators without any manual data pipeline. The
automated dropout alert mechanism, which fires in-app notifications when dropout
probability exceeds 65\%, exemplifies the platform's goal of shifting
institutional intervention from reactive to proactive.

The results validate that such integration is not only technically feasible
within a student-project scope but can be implemented to production-grade
standards: RS256 JWT authentication, immutable audit logs, Redis-backed rate
limiting, 13-container Docker orchestration, and a CI/CD pipeline covering all
three tiers of the stack.

% ─────────────────────────────────────────────────────────────────────────────
\section*{Acknowledgment}

The author sincerely thanks the faculty and staff of the Department of Computer
Science and Engineering, United Institute of Technology, Prayagraj, for their
guidance, feedback, and support throughout the development of EduIntel AI. The
department's computing infrastructure was instrumental in model training and
system deployment testing.

\textit{Prepared in partial fulfillment of Major Project requirements (Final
Year B.Tech), CSE Department, United Institute of Technology, Prayagraj.
Academic Year 2025--2026.}

% ─────────────────────────────────────────────────────────────────────────────
\begin{thebibliography}{00}

\bibitem{romero2010}
C. Romero and S. Ventura, ``Educational data mining: A review of the state of
the art,'' \textit{IEEE Trans. Syst. Man Cybern. C}, vol.~40, no.~6,
pp.~601--618, 2010.

\bibitem{kotsiantis2004}
S.~B. Kotsiantis, C. Pierrakeas, and P.~E. Pintelas, ``Predicting students'
performance in distance learning using machine learning techniques,''
\textit{Appl. Artif. Intell.}, vol.~18, no.~5, pp.~411--426, 2004.

\bibitem{mduma2019}
N. Mduma, K. Kalegele, and D. Machuve, ``A survey of machine learning
approaches and techniques for student dropout prediction,'' \textit{Data Sci.
J.}, vol.~18, no.~1, p.~14, 2019.

\bibitem{chen2016}
T. Chen and C. Guestrin, ``XGBoost: A scalable tree boosting system,'' in
\textit{Proc. ACM KDD}, 2016, pp.~785--794.

\bibitem{newman2021}
S. Newman, \textit{Building Microservices}, 2nd~ed. O'Reilly Media, 2021.

\bibitem{richardson2018}
C. Richardson, \textit{Microservices Patterns}. Manning Publications, 2018.

\bibitem{breiman2001}
L. Breiman, ``Random forests,'' \textit{Mach. Learn.}, vol.~45, no.~1,
pp.~5--32, 2001.

\bibitem{jones2015}
M. Jones, J. Bradley, and N. Sakimura, ``JSON Web Token (JWT),'' RFC~7519,
IETF, May 2015.

\bibitem{owasp2021}
OWASP Foundation, ``OWASP Top Ten 2021,'' 2021. [Online]. Available:
\url{https://owasp.org/www-project-top-ten/}

\bibitem{react}
React Team, ``React Documentation v18,'' 2023. [Online]. Available:
\url{https://react.dev/}

\bibitem{abramov}
S. Abramov, ``Redux: Predictable state container for JavaScript apps,'' Redux
Documentation, 2015. [Online]. Available: \url{https://redux.js.org/}

\bibitem{fastapi}
Tiangolo, ``FastAPI: Modern, fast web framework for building APIs with
Python,'' 2019. [Online]. Available: \url{https://fastapi.tiangolo.com/}

\bibitem{pedregosa2011}
F. Pedregosa et al., ``Scikit-learn: Machine learning in Python,'' \textit{J.
Mach. Learn. Res.}, vol.~12, pp.~2825--2830, 2011.

\bibitem{postgresql}
PostgreSQL Global Development Group, ``PostgreSQL 15 Documentation,'' 2023.
[Online]. Available: \url{https://www.postgresql.org/docs/15/}

\bibitem{docker}
Docker~Inc., ``Docker Documentation,'' 2024. [Online]. Available:
\url{https://docs.docker.com/}

\bibitem{guidotti2019}
R. Guidotti et al., ``A survey of methods for explaining black box models,''
\textit{ACM Comput. Surv.}, vol.~51, no.~5, pp.~1--42, 2019.

\end{thebibliography}

\end{document}
